\naslov{Galoisova teorija}

\podnaslov{Razširitve polj}

Pri določanju stopnje razširitve $F(a_1, \ldots, a_n)$ zapišemo vmesne
razširitve, in si pomagamo z naslednjo trditvijo:

\begin{trditev}
  Naj bo $F$ polje in $a_1, \ldots, a_n$ algebrajski nad $F$.
  Potem velja
  \begin{itemize}
  \item $\abs{F(a_1, \ldots, a_n) : F} = \abs{F(a_1, \ldots, a_n) : F(a_1,
	  \ldots, a_{n-1})} \cdots  \abs{F(a_1):F}$,
  \item $\abs{F(a_1, \ldots, a_n):F} \le \abs{F(a_1):F} \cdots \abs{F(a_n): F}$,
  \item če imata $\abs{F(a_1) : F}$ in $\abs{F(a_2) : F}$ največji skupni
	delitelj enak $1$, velja $\abs{F(a_1, a_2) : F} = \abs{F(a_1):F} \cdot
	\abs{F(a_2):F}$.
  \end{itemize}
\end{trditev}

Stopnje posameznih razširitev lahko ocenimo (najbolj enostavno, če $a_2 \notin
F(a_1)$, je $\abs{F(a_1, a_2):F(a_1)} \ge 2$), ali pa poiščemo minimalni
polinom, katerega stopnja je enaka stopnji razširitve.
Za minimalni polinom je dovolj, da je nerazcepen, in da ima ničlo $a_2$; pogosto
pa je že to težko.
Če je $p$ praštevilo, je $1 + x + x^2 + \cdots + x^{p-1}$ \pojem{ciklotoničen
  polinom}, ki je nerazcepen nad $\Q$.
Za polinome nad $\Q$ imaš na voljo tudi Eisensteinov kriterij:

\begin{trditev}
  Naj bo $p(x) = a_n x^n + \cdots + a_0$ polinom s koeficienti v $\Z$.
  Če obstaja praštevilo $p$, ki deli vse koeficiente razen $a_n$, in $p^2$ ne
  deli $a_0$, potem je $p$ nerazcepen nad $\Q$.
\end{trditev}

\begin{izrek}
  Naj bo $f \in F[X]$ nekonstanten nerazcepen polinom.
  Naslednje trditve so ekvivalentne:
  \begin{itemize}
  \item $f$ ima večkratno ničlo v razpadnem polju,
  \item $\operatorname{gcd}(f, f') \ne 1$,
  \item $f$ je polinom v $x^q$, kjer je $q$ karakteristika polja $F$,
  \item vse ničle v razpadnem polju so večkratne.
  \end{itemize}
\end{izrek}

\begin{trditev}
  Naj bo $f(X) \in F[X]$ polinom stopnje $n$.
  Potem $\abs{F(p):F}$ deli $n!$.
\end{trditev}

\podnaslov{Normalne in separabilne razširitve}

Za polinom $p(X) \in F[X]$ je razpadno polje $p$ najmanjša razširitev $E$, za
katero so vse ničle v $E$.
Določeno je do izomorfizma natančno, označimo ga z $F(p)$.
Če imamo izomorfizem polje $\sigma: F \to F'$, ga lahko vedno razširimo do
izomorfizma razpadnih polj.

\begin{definicija}
  Končna razširitev $E/F$ je \pojem{normalna}, če za vsak nerazcepen polinom
  $p(X) \in F[X]$ bodisi $p$ nima ničle v $E$, bodisi ima vse ničle v $E$.
\end{definicija}

\begin{izrek}
  Končna razširitev $E/F$ je normalna natanko tedaj, ko je $E$ razpadno polje
  nekega polinoma s koeficienti iz $F$.
\end{izrek}

\begin{definicija}
  Polinom $p(X) \in F[X]$ je \pojem{separabilen}, če ima same enostavne ničle.
  Končna razširitve $E/F$ je \pojem{separabilna}, če je za vsak $a \in E$
  minimalni polinom separabilen.
\end{definicija}

Če je karakteristika polja enaka $0$, in je polinom nerazcepen, so vse njegove
ničle enostavne.
V tem primeru je vsaka končna razširitev separabilna; polju s to lastnostjo
pravimo \pojem{perfektno}.
Poleg polj s karakteristiko $0$ vemo, da so tudi vsa končna polja perfektna.

\begin{izrek}[primitivni element]
  Vsaka separabilna razširitev je enostavna (tj.~obstaja $a \in E$, da je $E =
  F(a)$).
\end{izrek}

\podnaslov{Galoisove grupe in razširitve}

\begin{definicija}
  \pojem{Galoisova grupa} razširitve $E/F$ je grupa
  \[
	\operatorname{Gal}(E/F)
	= \{\sigma \in \Aut E \such \left. \sigma \right|_F = \id_F\}.
  \]
  Elementom te grupe pravimo \pojem{$F$-avtomorfizmi} polja $E$.
\end{definicija}

Galoisovo grupo polinoma $p(X) \in F[X]$ definiramo kot $\operatorname{Gal}(p) =
\operatorname{Gal}(F(p)/F)$.
Normalnim separabilnim razširitvam pravimo \pojem{Galoisove razširitve}.
Če ima $F$ karakteristiko $0$, je za polinom $p(X) \in F[X]$ je $F(p)/F$ vedno
Galoisova; to je končna razširitev, torej je normalna (ker je $F(p)$ razpadno
polje), in separabilna, ker ima polje karakteristiko $0$.
Pri določanju Galoisove grupe si pomagamo z naslednjo trditvijo:

\begin{trditev}
  Naj bo $E/F$ končna Galoisova razširitev.
  Potem je $\abs{\operatorname{Gal}(E/F)} = \abs{E:F}$.
  Če je $E/F$ končna separabilna, je $\abs{\operatorname{Gal}(E/F)} \le
  \abs{E:F}$.
\end{trditev}

Pogledamo torej stopnjo razširitve, da vemo, koliko elementov bo v grupi.
Potem te elemente poiščemo.
V primeru razpadnega polja pride prav naslednja lema:

\begin{lema}
  Če je $a \in E$ ničla polinoma $p(X) \in F[X]$ ter $\sigma \in
  \operatorname{Gal}(E/F)$, je tudi $\sigma(a)$ ničla $p(X)$.
\end{lema}

Avtomorfizme iščemo tako, da si izberemo, kam bomo slikali ničle.
Če jih najdemo dovolj, vemo, da je to to.
Potem moramo le določiti, kaj točno je grupa.
Najbolj enostavno je preveriti komutativnost in red elementov (torej poiskanih
avtomorfizmov).
Spomnimo se klasifikacijo končnih Abelovih grup:

\begin{izrek}
  Vsaka končna Abelova grupa je vsota cikličnih grup, katerih red je potenca
  praštevila.
\end{izrek}

Rezultat, ki bizarno redko pride prav, je fundamentalni izrek Galoisove teorije:

\begin{izrek}
  Naj bo $E/F$ končna Galoisova razširitev.
  \begin{itemize}
  \item Predpisa $L \mapsto \operatorname{Gal}(E/L)$ in
	\[
	  G \mapsto \{x \in E \such g \cdot x = x \;\forall g \in G\}
	\]
	sta paroma inverzni preslikavi med vmesnimi polji $F \subseteq L \subseteq
	E$ in podgrupami Galoisove grupe.
  \item Za poljubni vmesni polji $F \subseteq L \subseteq M \subseteq E$ velja
	$\abs{M:L} = \abs{\operatorname{Gal}(E/L) : \operatorname{Gal}(E/M)}$.
  \item Za vmesno polje $F \subseteq L \subseteq E$ je $L/F$ normalna natanko
	tedaj, ko je $\operatorname{Gal}(E/L) \edinka \operatorname{Gal}(E/F)$.
  \end{itemize}
\end{izrek}

\podnaslov{Rešljivost grup}

\begin{definicija}
  Grupa $G$ je \pojem{rešljiva}, če obstaja končno zaporedje podgrup
  \[
	\{1\} = G_0 \le G_1 \le \cdots \le G_k = G,
  \]
  za katere velja $G_i \edinka G_{i+1}$ in $G_{i+1} / G_i$ Abelova.
\end{definicija}

Naslednje grupe so dokazano rešljive:
\begin{itemize}
\item Abelove grupe,
\item končne $p$-grupe,
\item $D_{2n}$,
\item vsaka grupa reda $pq$, kjer sta $p,q$ praštevili
\end{itemize}
Vse podgrupe in kvocienti rešljivih grup so rešljivi.
Če sta $N$ in $G/N$ rešljivi za edinko $N \edinka G$, je tudi $G$ rešljiva.

\naslov{Moduli}

Modul je algebrajska struktura nad kolobarjem $R$, v kateri dovolimo seštevanje
in množenje s skalarjem z leve.
Veljajo vse običajne definicije, kar se tiče algebrajskih struktur.
Če je $X = \{x_1, \ldots, x_n\}$, je podmodul, generiran z $X$, enak
\[
  \sk{X} = \{r_1 x_1 + \cdots + r_n x_n \such r_i \in R\}.
\]
Če je modul generiran z enim elementom, pravimo, da je \pojem{cikličen}.
Vsak enostaven modul (tj.~tak, ki nima netrivialnih podmodulov) je cikličen.

Množici $X$, za katero je $M = \sk{X}$ in za katero iz enakosti $r_1 x_1 +
\cdots + r_n x_n = 0$ sledi $r_i = 0$ za vsak $i$, pravimo \pojem{baza}.
Če ima modul bazo, je \pojem{prost}.
Velja naslednja karakterizacija:

\begin{izrek}
  Naj bo $M$ $R$-modul.
  Naslednje trditve so ekvivalentne:
  \begin{itemize}
  \item $M$ je prost $R$-modul
  \item Obstaja indeksna množica $\Lambda$, da je $M$ kot $R$-modul izomorfen
	\[
	  M \cong \bigoplus_{\lambda \in \Lambda} R
	\]
  \item Obstaja množica $X$ in preslikava $\iota: X \to M$, da za vsak $R$-modul
	$N$ in vsako preslikavo $\kappa: X \to N$ obstaja natanko en homomorfizem
	$f: M \to N$, za katerega je $f \circ \iota = \kappa$.
  \end{itemize}
\end{izrek}

Zadnji točki pravimo \pojem{univerzalna lastnost}.
Pripada ji naslednji komutativni diagram:
\[\begin{tikzcd}
	X & N \\
	M
	\arrow["\kappa", from=1-1, to=1-2]
	\arrow["\iota"', from=1-1, to=2-1]
	\arrow["f"', from=2-1, to=1-2]
  \end{tikzcd}\]

% LocalWords:  ciklotoničen separabilne Galoisove Galoisovo separabilnim
